\chapter*{Tóm tắt Khóa luận}
\label{abstract}

\section*{Tiếng Việt}

Đề tài \textbf{"Đề xuất khung phân tích hỗ trợ sinh viên phân tích dữ liệu tuyển dụng"} được thực hiện nhằm giải quyết bài toán thiếu hụt công cụ hỗ trợ sinh viên đánh giá định lượng khoảng cách năng lực của bản thân so với yêu cầu thực tế của thị trường lao động. Trong bối cảnh dữ liệu tuyển dụng trực tuyến ngày càng phong phú nhưng thiếu tính chuẩn hóa, các hệ thống đối soát hồ sơ truyền thống chủ yếu dựa trên so khớp từ khóa chính xác (Exact Keyword Matching) dẫn đến nhiều hạn chế về tính chính xác và tính linh hoạt ngữ nghĩa. Luận văn đề xuất và hiện thực hóa hệ thống \textit{CareerNova} --- một nền tảng thông minh tích hợp Pipeline Dữ liệu lớn (ETL), Mô hình ngôn ngữ lớn (LLM) và Xử lý ngôn ngữ tự nhiên (NLP) nhằm trực quan hóa xu hướng thị trường và tự động hóa quy trình đối soát hồ sơ (CV Matching).

Hệ thống được thiết kế với ba phân hệ cốt lõi: (1) Phân hệ thu thập và chuẩn hóa dữ liệu tự động từ các nền tảng tuyển dụng lớn, áp dụng kỹ thuật \textit{Structured Outputs} từ LLM kết hợp thuật toán chống ảo giác để đảm bảo chất lượng dữ liệu đầu vào; (2) Phân hệ trực quan hóa dữ liệu thị trường với các Dashboard tương tác giúp người dùng nắm bắt xu hướng kỹ năng, quyền lợi và nhu cầu tuyển dụng theo thời gian; và (3) Phân hệ đối soát ngữ nghĩa dựa trên vector nhúng SBERT (Sentence-BERT) và trọng số TF-IDF được tính toán động từ thị trường, vượt trội hơn so với các phương pháp so khớp từ khóa truyền thống.

Kết quả thực nghiệm trên môi trường kiểm thử cho thấy hệ thống hoạt động ổn định và chính xác. Pipeline ETL đã thu thập và làm sạch thành công 42.520 tin tuyển dụng, đạt tỷ lệ dữ liệu sạch 74,66\%. Trong bài toán trích xuất kỹ năng bằng LLM, hệ thống đạt độ chính xác (Precision) 87,5\%, độ phủ (Recall) 84,2\% và F1-Score 85,8\%. Toàn bộ 15/15 kịch bản kiểm thử tích hợp (Integration Test) đều vượt qua thành công, với thời gian phản hồi trung bình của API đối soát đạt 1,45 giây. Luận văn đóng góp một giải pháp có tính ứng dụng thực tiễn cao, cân bằng giữa chất lượng mô hình LLM và hiệu năng của hệ thống vector cục bộ, tạo ra một công cụ hỗ trợ ra quyết định nghề nghiệp đắc lực cho sinh viên trong thời đại số.

\section*{Abstract (English)}

This thesis, titled \textbf{"Proposal of an Analytical Framework to Assist Students in Analyzing Recruitment Data"}, addresses a critical gap in career development support: the absence of a quantitative tool helping students accurately assess their skill gaps against the real-world demands of the IT labor market. As online recruitment data grows increasingly abundant yet unstructured, conventional profile-matching systems that rely on Exact Keyword Matching suffer from significant limitations in semantic flexibility and precision. To overcome these challenges, this thesis proposes and implements the \textit{CareerNova} system --- an intelligent platform integrating Big Data ETL Pipelines, Large Language Models (LLMs), and Natural Language Processing (NLP) to visualize market trends and automate the CV Matching process.

The system is designed with three core subsystems: (1) An automated data collection and normalization pipeline from major recruitment platforms, applying \textit{Structured Output} techniques from LLMs combined with an anti-hallucination validation mechanism to ensure high-quality input data; (2) An interactive data visualization module with dynamic dashboards enabling users to explore trends in skills, benefits, and hiring demands over time; and (3) A semantic matching engine based on SBERT (Sentence-BERT) embeddings and dynamically computed TF-IDF market weights, significantly outperforming traditional keyword-based approaches.

Experimental results demonstrate stable and highly accurate system performance. The ETL pipeline successfully collected and cleaned 42,520 job postings, achieving a clean data rate of 74.66\%. In the LLM-based skill extraction task, the system achieved a Precision of 87.5\%, Recall of 84.2\%, and F1-Score of 85.8\%. All 15/15 integration test cases passed successfully, with an average API response time of 1.45 seconds. This thesis contributes a highly practical solution that balances the expressiveness of LLM-based extraction with the performance of local vector search, ultimately creating a data-driven career decision support tool for students in the digital era.
